\chapter{Event Sources}

The previous chapter explained various frameworks for supporting event tracing.
This chapter discusses some of the most common types of events used in Linux,
focusing on their registration, implementation and interaction with the
aforementioned frameworks. In general, an event is triggered when internal or
external conditions are met, causing the system to execute logic that reacts to
the occurrence. The lifetime of an event is managed by a tracing framework, and
the event may use one or more frameworks to produce the tracing data. A
distinction is made between \textit{static} events and \textit{dynamic} events:
static events are placed by kernel developers in pre-determined locations and
statically compiled in the kernel, while dynamic events are registered during
runtime and can be registered in non-predetermined locations.

\section{Tracepoints}

A \textit{tracepoint} is a system that calls an arbitrary function when
triggered, which is referred to as a "callback". Trigger points are statically
placed all over the kernel's codebase between regular code; there are more than
1000 tracepoints available in Linux 7.2. A tracepoint can be either set as "on",
meaning a callback is connected to it, or "off" meaning no callback is present.
The callback can accept custom parameters for tracing purposes. The usage of
tracepoints is conceptually simple, which will be explained with several
examples.

Consider the following code:

\begin{Verbatim}[numbers=left, numbersep=5pt, frame=single]
void foo(void) {
  int x = bar();
  /* ... */
  trace_foo(x); /* Tracepoint trigger placed here */
  /* ... */
  return baz(x);
}
\end{Verbatim}

Here the function \verb|foo()| calls at some point \verb|trace_foo()| on line
$4$ with the variable \verb|x| passed as an argument. \verb|trace_foo()| is a
function generated by the tracpoint API for a particular tracepoint called
"foo", along side with other functions. In this case, if the tracepoint is
currently active, then this function will jump to callback function specified
when registering the tracepoint. The callback can use the \verb|x| variable for
any purposes, such as to provide it to the user by writing it to the ftrace or
the perf ring buffer.

\subsection{Tracepoint related types}

The key types that the tracepoint API defines are \verb|struct tracpoint| which
contains all information related to the tracepoint, \verb|struct tracepoint_func|
containing the callback, and \path{struct tracepoint_ext}
specifying extended attributes. They are reported below from the file
\path{$LINUX/include/linux/tracepoint-defs.h:26}, where \verb|$LINUX|
represents the root of the Linux kernel source code:

\begin{Verbatim}[numbers=left, numbersep=5pt, frame=single]
struct tracepoint_func {
  void *func;
  void *data; int prio;
};

struct tracepoint_ext {
  int (*regfunc)(void);
  void (*unregfunc)(void);
  /* Flags. */
  unsigned int faultable:1;
};

struct tracepoint {
  const char *name;    /* Tracepoint name */
  struct static_key_false key;
  struct static_call_key *static_call_key;
  void *static_call_tramp;
  void *iterator;
  void *probestub;
  struct tracepoint_func __rcu *funcs;
  struct tracepoint_ext *ext;
};
\end{Verbatim}

The most relevant members of \verb|struct tracepoint| are its "name" on line
$14$, which is used by the user to identify a particular tracepoint, and the
"funcs" member on line $20$ used to specify one or more functions to be called
when the tracepoint is executed (which are of type \verb|struct tracepoint_func|).
The "ext" (extended) member on line $21$ contains optional
functions to be called when registering or unregistering the tracepoint.
Additionally, the static-key and static-call related fields are performance
optimizations which will be briefly explained later in this section.

\subsection{Tracepoint creation}

To create a tracepoint, the macro \verb|DECLARE_TRACE()| should be used. This
macro takes three arguments: the name of the tracepoint, the prototype of the
callback, and the name of the arguments what should be passed to the callback
(relevant to the specific scope where the tracepoint is placed). When expanded,
it sets up the necessary static storage for the tracepoint and generates a set
of functions to interact with it.

For example, consider the following macro statement:

\begin{Verbatim}[numbers=left, numbersep=5pt, frame=single]
DECLARE_TRACE(foo_tp,
  TP_PROTO(bool is_bar),
  TP_ARGS(bar));
\end{Verbatim}

This declares a tracepoint called \verb|foo_tp| which calls a function that
looks like \path{foo_function(bool is_bar)}. This function will be called with
the \verb|bar| argument, like \verb|foo_function(bar)|.

Consider now a real example from \path{$LINUX/include/trace/events/sched.h}:

\begin{Verbatim}[numbers=left, numbersep=5pt, frame=single]
#undef TRACE_SYSTEM
#define TRACE_SYSTEM sched

#if !defined(_TRACE_SCHED_H) || defined(TRACE_HEADER_MULTI_READ)
#define _TRACE_SCHED_H

#include <linux/tracepoint.h>

DECLARE_TRACE(sched_entry,
  TP_PROTO(bool preempt),
  TP_ARGS(preempt));

#endif /* _TRACE_SCHED_H */

/* This part must be outside protection */
#include <trace/define_trace.h>
\end{Verbatim}

In the first line, the \verb|TRACE_SYSTEM| symbol specifies the subdirectory of
tracefs where this tracepoint will be registered for access to the user. The
code follows by including the file \path{$LINUX/tracepoint.h} in line $7$
which contains the \verb|DECLARE_TRACE()| macro definition, and proceeds to
declare the \verb|sched_entry| tracepoint on line $9$. It finally includes
\path{trace/define\_trace.h} on line $16$ which defines the various functions
for the tracepoint and static storage, such as functions for registration,
de-registration, calling the callback, and a function to check if the tracepoint
is active, among others. After all macro expansions, the following functions are
generated and can be used throughout the codebase:

\begin{itemize}

  \item \verb|register_trace_sched_entry()|: takes the callback and sets it up
    to be called by the tracepoint.

  \item \verb|unregister_trace_sched_entry()|: unregisters the tracepoint.

  \item \verb|trace_call__sched_entry()|: calls the callback.

  \item \verb|trace_sched_entry()|: calls the callback if the tracepoint is.
    enabled.

  \item \verb|trace_sched_entry_enabled()|: checks if the tracepoint is enabled.

\end{itemize}

It should be noted that the \verb|define_trace.h| header needs to be placed
outside any inclusion guards so that it will be always included by the
preprocessor. This is used for a trick with the preprocessor where the
\verb|define_trace.h| header includes the file that is including it, therefore
including itself again and redefining the macro \verb|DECLARE_TRACE()|, but
using \verb|#define| guards to change its definition. For example, if file "X"
includes \verb|define_trace.h|, then it will include itself, "X", which will
include \verb|define_trace.h| another time, and so on; this cycle may repeat
multiple times, until some condition is met. The reason is to redefine the
\verb|DECLARE_TRACE()| macro in different ways to generate different code. This
is also used in other parts of the kernel, namely "event tracing" which will be
explained later in this section.

\subsection{static\_call and static\_key mechanism}

After declaring a tracepoint with \verb|DECLARE_TRACE()|, the function
\verb|trace_sched_entry()| can be used to trigger the tracepoint. The callstack
looks like the following, where the function on top calls the function
underneath:

\begin{verbatim}
trace_sched_entry(...)
  __do_trace_sched_entry(...)
      static_call(tp_func_sched_entry)
\end{verbatim}

The \textit{static call} on the last line is a mechanism that uses code patching
to replace global function pointers. It allows directs calls to be used instead
of indirect calls. In this particular case, instead of having a pointer to the
callback function stored somewhere in the tracepoint structure which needs to be
loaded from memory, a call to its address is directly patched where the
\verb|static_call()| call lives therefore removing one layer of indirection.
There may be multiple callbacks registered, in this case the call will be
patched with an iterator which simply loops over all registered callbacks and
calls them.

Before calling \verb|static_call()|, a check is made to determine if the
tracepoint is active or not by checking if the callback address is present. To
avoid a conditional branch instruction in the if statement, the code is patched
at runtime via a mechanism called \verb|static_key| that replaces the if
statement with a jump or a NOP instruction depending on whether the event is
active or not. Note that this is different than \verb|static_call()| because it
replaces a conditional branch (compare + jump instructions) and not a function
call. This decreases the runtime performance impact of a disabled tracepoint, as
the only effective penalty is overall increased code section size due to the
tracepoint generated functions.

\section{Event Tracing}

A tracepoint provides a solution to call a callback function when triggered.
Additional work is necessary to make this system interact with the rest of the
tracing frameworks in the kernel, such as ftrace and perf. To simplify this, the
macro \verb|TRACE_EVENT()| should be used instead of \verb|DEFINE_TRACE()|.
Creating a tracepoint using this macro is referred to as "event tracing", which
is not to be confused with "perf events".

The \verb|DEFINE_TRACE()| macro is responsible for creating a proper tracepoint,
creating a callback that can be hooked to the tracepoint recording some
formatted data into the tracer ring buffer, supporting both perf and ftrace. To
accomplish this, the \verb|TRACE_EVENT()| macro is broken into six components,
which correspond to the parameters of the macro:

\begin{itemize}
  \item \textit{name}: the name of the tracepoint to be created

  \item \textit{prototype}: the prototype for the tracepoint callbacks

  \item \textit{args}: the arguments that match the prototype

  \item \textit{struct}: the structure that a tracer could use (but is not
  required to) to store the data passed into the tracepoint

  \item \textit{assign}: code written in C to assign the data to the structure.

  \item \textit{print}: code written in C to output the structure in human
  readable ASCII format.

\end{itemize}

For example, consider the trace event in
\path{$LINUX/fs/fuse/fuse_trace.h:126}:

\begin{Verbatim}[numbers=left, numbersep=5pt, frame=single]
TRACE_EVENT(fuse_request_end,
  TP_PROTO(const struct fuse_req *req),

  TP_ARGS(req),

  TP_STRUCT__entry(
    __field(dev_t,    connection)
    __field(uint64_t,  unique)
    __field(uint32_t,  len)
    __field(int32_t,  error)
  ),

  TP_fast_assign(
    __entry->connection  =  req->chan->conn->dev;
    __entry->unique    =  req->in.h.unique;
    __entry->len    =  req->out.h.len;
    __entry->error    =  req->out.h.error;
  ),

  TP_printk("connection %u req %llu len %u error %d", __entry->connection,
      __entry->unique, __entry->len, __entry->error)
);
\end{Verbatim}

This code creates the trace event \verb|trace_fuse_request_end| using the
specified parameters to format the event. The expansion of this macro uses the
same trick as \verb|DEFINE_TRACE()| by redefining itself a total of 7 times
which will not be covered here due to the high complexity of the expansion.
After the expansion, the function \verb|trace_fuse_request_end()| can be called
anywhere in the code to trigger the tracepoint, which will use ftrace to produce
its output, formatted as specified in \verb|TP_printk()| on line $20$ using the
variables from \verb|TP_fast_assign()| on line $13$. Additionally, the
\verb|TRACE_EVENT()| macro registers entries in tracefs to disable and enable
the event (\path{$TRACEFS/events/fuse/fuse_request_end/enable}).

\subsection{Event classes}

Since \verb|TRACE_EVENT()| generates a lot of code and increases the size of the
kernel binary, the macro \verb|DECLARE_EVENT_CLASS()| and \verb|DEFINE_EVENT()|
should be used for tracepoints that share some code. The class defines the set
of variables and formatting which can be reused by multiple events belonging to
that class through \verb|DEFINE_EVENT()|.

\section{Syscall Tracepoints}

For various use cases, it may be necessary to trace accesses to systemcalls. For
this reason, special tracepoints are placed specifically for systemcalls which
are triggered when the function is entered and exited. They are used to trace
both systemcall name and its arguments to a tracing framework, for any
systemcall. An explaination of the implementation is now provided.

A systemcall is defined by the \verb|SYSCALL_DEFINEX()| family of macros. Taking
\path{SYSCALL_DEFINE0()} as an example, the macro expands to the following code
on x86, extracted from
\path{$LINUX/arch/x86/include/asm/syscall_wrapper.h:244}:

\begin{Verbatim}[numbers=left, numbersep=5pt, frame=single]
#define SYSCALL_DEFINE0(sname)            \
  SYSCALL_METADATA(_##sname, 0);          \
  static long __do_sys_##sname(const struct pt_regs *__unused);  \
  __X64_SYS_STUB0(sname)            \
  __IA32_SYS_STUB0(sname)            \
  static long __do_sys_##sname(const struct pt_regs *__unused)
\end{Verbatim}

Notice on line 2 the macro \verb|SYSCALL_METADATA|, which expands to the
following code, extracted from \path{\$LINUX/include/linux/syscalls.h:178}:

\begin{Verbatim}[numbers=left, numbersep=5pt, frame=single]
#define SYSCALL_METADATA(sname, nb, ...)      \
  static const char *types_##sname[] = {      \
    __MAP(nb,__SC_STR_TDECL,__VA_ARGS__)    \
  };              \
  static const char *args_##sname[] = {      \
    __MAP(nb,__SC_STR_ADECL,__VA_ARGS__)    \
  };              \
  SYSCALL_TRACE_ENTER_EVENT(sname);      \
  SYSCALL_TRACE_EXIT_EVENT(sname);      \
  static struct syscall_metadata __used      \
    __syscall_meta_##sname = {        \
    .name     = "sys"#sname,      \
    .syscall_nr  = -1,  /* Filled in at boot */  \
    .nb_args   = nb,        \
    .types    = nb ? types_##sname : NULL,  \
    .args    = nb ? args_##sname : NULL,  \
    .enter_event  = &event_enter_##sname,    \
    .exit_event  = &event_exit_##sname,    \
    .enter_fields  = LIST_HEAD_INIT(__syscall_meta_##sname.enter_fields), \
  };              \
  static struct syscall_metadata __used      \
    __section("__syscalls_metadata")      \
   *__p_syscall_meta_##sname = &__syscall_meta_##sname;
\end{Verbatim}

On line 8 and 9 the macros \path{SYSCALL_TRACE_ENTER_EVENT()} and
\path{SYSCALL_TRACE_EXIT_EVENT()} are responsible for creating the tracepoint,
and on line 10 a \verb|struct syscall_metadata| is created which is used to save
some information about the syscall and its arguments. Finally on line 21 a
pointer to the metadata struct is created in the section
\verb|__syscalls_metadata| which contains all metadata from all systemcalls.

Whenever a systemcall is executed by the user, the following functions are
called from top to bottom:

\begin{verbatim}
entry_SYSCALL_64
   do_syscall_64
      syscall_enter_from_user_mode
         syscall_trace_enter
            trace_syscall_enter
              trace_sys_enter
\end{verbatim}

The last function, \verb|trace_sys_enter|, is a trace event defined in
\path{$LINUX/include/trace/events/syscalls.h}:

\begin{Verbatim}[numbers=left, numbersep=5pt, frame=single]
TRACE_EVENT_SYSCALL(sys_enter,

  TP_PROTO(struct pt_regs *regs, long id),

  TP_ARGS(regs, id),

  TP_STRUCT__entry(
    __field(  long,    id    )
    __array(  unsigned long,  args,  6  )
  ),

  TP_fast_assign(
    __entry->id  = id;
    syscall_get_arguments(current, regs, __entry->args);
  ),

  TP_printk("NR %ld (%lx, %lx, %lx, %lx, %lx, %lx)",
      __entry->id,
      __entry->args[0], __entry->args[1], __entry->args[2],
      __entry->args[3], __entry->args[4], __entry->args[5]),

  syscall_regfunc, syscall_unregfunc
);
\end{Verbatim}

This last macro creates the actual tracepoint that logs all systemcalls, using
the metadata registered from \verb|SYSCALL_DEFINEX()| via
\verb|SYSCALL_METADATA()| for printing. This also creates an entry in tracefs
for each systemcall to enable tracing for that specific systemcall
(\$TRACEFS/events/syscalls/\$SYSCALL).

%\section{User Statically-Defined Tracing}
%TODO

\section{Kprobes and Kretprobes}

\textit{Kprobes} are events that can be dynamically placed almost anywhere on
the kernel. They use code patching to replace any instruction with a breakpoint
instruction; when triggered, it calls one or more "probe" functions which may
perform some tracing operations, and restores the registers when they complete.
The instruction that was replaced will be single-stepped (executed or emulated)
by the probe, so that the program will continue to run correctly as originally
expected. Such instrumentation does not require any special infrastructure on
the traced function, which makes kprobes particularly useful and versatile. The
kprobe infrastructure is available with the kernel compile option
\verb|CONFIG_KPROBES|.

Two different kinds of probes exist: \textit{kprobes} which are often placed at
the beginning of functions, and \textit{kretprobes} which are called at the end
of the function and may alter the return value. A kprobe or kretprobe sets up a
call to two functions, if specified: a \textit{pre-handler}, which is called
before single-stepping the replaced instruction, and a \textit{post-handler}
which is instead executed after. Figure \ref{fig:execution_of_a_kprobe} shows
these phases in time. Additionally, kprobes can be attached to both regular code
and to ftrace's function tracing inftrastructure to be attached to each
function.

\begin{figure}[htbp]
\centering
\begin{tikzpicture}[
    scale=1.0,
    every node/.style={
        transform shape
    },
    node distance=1.5cm,
    block/.style={
      rectangle,
      draw=gray!80,
      fill=gray!5,
      rounded corners=4px,
      minimum width=3.2cm,
      minimum height=1.2cm,
      font=\small\bfseries,
      align=center,
      line width=1.2pt
    },
    flow/.style={
      ->,
      line width=1.2pt,
      draw=gray!80
    }
]

    % =========================================================================
    % FLOWCHART NODES (Top Row)
    % =========================================================================
    \node[block, text width=2.8cm] (int3) {int3 hit\\[0.35em]\mdseries\footnotesize\texttt{do\_int3}};

    \node[block, right=1.8cm of int3, text width=3.8cm] (prehandler) {Call pre\_handler\\[0.2em]Set up single-stepping};

    \node[block, right=1.8cm of prehandler, text width=2.8cm] (singlestep)
    {single-step\\probed instruction};

    % =========================================================================
    % FLOWCHART NODES (Bottom Row)
    % =========================================================================
    \node[block, below=2.2cm of singlestep, text width=3.8cm] (posthandler) {Call post\_handler\\[0.2em]Resume Execution};

    \node[block, left=1.8cm of posthandler, text width=2.8cm] (nextinst) {next instruction\\[0.2em]after probe};

    % =========================================================================
    % CONNECTORS
    % =========================================================================
    \draw[flow] (int3) -- (prehandler);
    \draw[flow] (prehandler) -- (singlestep);

    % Path dropping down from single-step to post-handler with label
    \draw[flow] (singlestep.south) -- (posthandler.north);

    \draw[flow] (posthandler) -- (nextinst);

    % =========================================================================
    % SIDE LABEL
    % =========================================================================
    %\node[font=\bfseries\small, rotate=90] at ($(int3.west) + (-0.6cm, 0)$) {probe};

\end{tikzpicture}
\caption{Key phases furing the execution a Kprobe.}
\label{fig:execution_of_a_kprobe}
\end{figure}


\subsection{Kprobe interface}

A kprobe is represented by the type \verb|struct kprobe| which is reported
below, with comments explaining the various members:

\begin{verbatim}
typedef int (*kprobe_pre_handler_t) (struct kprobe *, struct pt_regs *);
typedef void (*kprobe_post_handler_t) (struct kprobe *, struct pt_regs *,
              unsigned long flags);

struct kprobe {
  struct hlist_node hlist;
  /* list of kprobes for multi-handler support */
  struct list_head list;
  /*count the number of times this probe was temporarily disarmed */
  unsigned long nmissed;
  /* location of the probe point */
  kprobe_opcode_t *addr;
  /* Allow user to indicate symbol name of the probe point */
  const char *symbol_name;
  /* Offset into the symbol */
  unsigned int offset;
  /* Called before addr is executed. */
  kprobe_pre_handler_t pre_handler;
  /* Called after addr is executed, unless... */
  kprobe_post_handler_t post_handler;
  /* Saved opcode (which has been replaced with breakpoint) */
  kprobe_opcode_t opcode;
  /* copy of the original instruction */
  struct arch_specific_insn ainsn;
  /*
   * Indicates various status flags.
   * Protected by kprobe_mutex after this kprobe is registered.
   */
  u32 flags;
};
\end{verbatim}

A kprobe can be registered and unregistered with the functions
\verb|register_kprobe()| and \path{unregister_kprobe()}. Kprobes can be
temporarely enabled with \verb|enable_kprobe()| and disabled with
\verb|disable_kprobe()| after registration. The list of registered kprobes is
exposed to the user by reading the file \path{$DEBUGFS/kprobes/list} where
\verb|$DEBUGFS| is the location where the \textit{debugfs} filesystem is
mounted.

To showcase and explain the kprobe interface, an example kernel module will be
analyzed. The module register a kprobe on the \verb|do_sys_openat2()| function
which is called by the \verb|openat2| systemcall, and traces the name of the
file that the user wants to open:

\begin{Verbatim}[numbers=left, numbersep=5pt, frame=single]
#include <linux/module.h>
#include <linux/kprobes.h>

#define MAX_FILENAME_LEN 128

/* This gets executed before single-stepping the probed instruction */
static int hello_kprobe_pre_handler(struct kprobe *p, struct pt_regs *regs)
{
  const char __user *user_filename =
    (const char __user *)regs_get_kernel_argument(regs, 1);
  if (!user_filename)
    return 0;

  char *filename = kmalloc(MAX_FILENAME_LEN, GFP_ATOMIC);
  if (!filename)
    return 0;

  if (strncpy_from_user(filename, user_filename, MAX_FILENAME_LEN-1) > 0) {
    filename[MAX_FILENAME_LEN-1] = '\0';
    printk(KERN_INFO "Kprobe: openat -> %s\n", filename);
  }

  kfree(filename);
  return 0;
}

/* This gets executed after single-stepping the probed instruction */
static void hello_kprobe_post_handler(struct kprobe *kprobe,
                                      struct pt_regs *regs,
                                      unsigned long flags)
{ /* does nothing */ }

struct kprobe kp = (struct kprobe) {
  .symbol_name  = "do_sys_openat2",
  .pre_handler  = hello_kprobe_pre_handler,
  .post_handler = hello_kprobe_post_handler,
};

static int __init hello_kprobe_init(void)
{
  printk(KERN_INFO "Kprobe: Hello, kprobe!\n");

  int ret = register_kprobe(&kp);
  if (ret < 0) {
    printk(KERN_ERR "Kprobe: error registering kprobe\n");
    return -1;
  }
  return 0;
}

static void __exit hello_kprobe_cleanup(void)
{
  unregister_kprobe(&kp);
  printk(KERN_INFO "Kprobe: unregistered");
}

module_init(hello_kprobe_init);
module_exit(hello_kprobe_cleanup);

MODULE_AUTHOR("Giovanni Santini");
MODULE_DESCRIPTION("kprobe example");
MODULE_LICENSE("GPL v2");
\end{Verbatim}

On line 7 and line 28 the kprobe's pre-handler \verb|hello_kprobe_pre_handler()|
and post-handler \verb|hello_kprobe_post_handler()| functions are defined. Both
pre and post handlers can access the CPU registers via the struct \verb|struct pt_regs|.
Looking at the pre-handler, the values of the registers correspond to
the argument to the function according to the C ABI for the specific CPU
architecture used, since the kprobe was placed in a function. The pre-handler
gets the second argument from the function using
\verb|regs_get_kernel_argument()| on line 10, which corresponds to the filename,
then copies its value from userspace to kernelspace on line 18 and logs it on
line 20 using \verb|printk()| which writes it to a special ringbuffer used for
kernel prints which is accessible from \path{$PROCFS/kmsg}. On line 33 the
code defines the kprobe \verb|kp| with its pre and post handler, and the
location where it should be placed. The specific address of the kprobe will be
resolved from the symbol specified in the \verb|symbol_name| member. Finally on
line 43 the kprobe is registered when the module is loaded, and unregistered on
line 53 when the module is unloaded.

\subsection{Kprobe implementation}

Most of the implementation of the kprobe interface is located in the file
\path{$LINUX/kernel/kprobes.c}. Each registered kprobe is stored in the hash
table \verb|kprobe_table| which is indexed by the address where it is placed.
Insertion and deletion to this table is guarded by the mutex \verb|kprobe_mutex|
to prevent such operations from executing simultaneously. A list of blacklisted
address ranges \verb|kprobe_blacklist| is used for areas where kprobes are not
allowed to be registered.

An example that showcases the use of \verb|kprobe_table| is the function
\verb|get_kprobe()| present in \path{$LINUX/kernel/kprobes.c:378} which
returns the \verb|struct kprobe| from an address where a kprobe was placed:

\begin{Verbatim}[numbers=left, numbersep=5pt, frame=single]
struct kprobe *get_kprobe(void *addr)
{
  struct hlist_head *head;
  struct kprobe *p;

  head = &kprobe_table[hash_ptr(addr, KPROBE_HASH_BITS)];
  hlist_for_each_entry_rcu(p, head, hlist,
         lockdep_is_held(&kprobe_mutex)) {
    if (p->addr == addr)
      return p;
  }

  return NULL;
}
\end{Verbatim}

In the code above, a list of kprobes is retrieved from the hashmap in constant
time complexity by using an hash of the pointer where the kprobe is placed.
Then, for each kprobe in the list, if its address actually corresponds to the
query then it is returned. Otherwise the function returns a \verb|NULL| pointer.

When registering a kprobe regularly (not through ftrace),
\verb|register_kprobe()| first retrieves the address of the probed instruction
from the symbol, if specified, via \verb|kallsyms_lookup_name()|. This performs
a binary search of the symbol name from a compressed array containing a mapping
from each symbol to its address. The registration continues by calling
\verb|prepare_kprobe()| which sets up internal structures, it then adds the hash
of the kprobe to the \verb|kprobe_table| map and calls \verb|arm_kprobe()| only
if the kprobe was not set to be disabled by default. This last function is
responsible to replace the instrumented instruction in memory with a breakpoint
call (\verb|int3| on x86). It finally sets up the emulation of the probed
instruction, which emulate it instead of calling it directly, avoiding the need
for a second breakpoint to call the post-handler. After registering, if the
kernel is compiled with \verb|CONFIG_OPTPROBES|, it tries to optimize the break
instruction with a jump. This does not work if the kprobe was based on ftrace.
The optimizer logic is quite complex and will not be covered here.

The \verb|int3| breakpoint handler for x86 is defined in
\path{$LINUX/arch/x86/kernel/trace.c}. It calls \verb|do_int3()| which
interfaces with the kprobe infrastructure to call the registered kprobe, if
present, as shown in the following extract:

\begin{Verbatim}[numbers=left, numbersep=5pt, frame=single]
static bool do_int3(struct pt_regs *regs)
{
  int res;
  /* ... */
#ifdef CONFIG_KPROBES
  if (kprobe_int3_handler(regs))
    return true;
#endif
  /* ... */
  return res == NOTIFY_STOP;
}
NOKPROBE_SYMBOL(do_int3);
\end{Verbatim}

The \verb|kprobe_int3_handler()| called on line 6 uses two per-cpu variable
called \verb|current_kprobe| and \verb|kprobe_ctrlblk| to implement the
following logic:

\begin{itemize}

  \item If \verb|current_kprobe| is empty, meaning that a kprobe is not running,
    then it sets \path{current_kprobe}. If specified, it calls the
    \verb|kp->pre_handler| and proceeds to emulate the probed instruction and
    call the post handler, or to setup the CPU's instruction pointer to the
    original instruction. In this case, the instruction pointer will be set when
    the trap returns, the instruction will be executed, and a new trap will be
    called with \verb|int3| and with \verb|current_kprobe| not empty. This is
    necessary because jump instructions cannot be emulated or executed during
    the single step since they would.

  \item if \verb|current_kprobe| is non empty, which happens after single
    stepping the instruction and entering the trap again from the previous
    point. This returns to the normal execution by setting the instruction
    pointer to the original place (just after the probed instruction).

\end{itemize}

A kretprobe allows the user do register a callback that gets executed when the
function returns. To implement this, the kretprobe registers a kprobe
\verb|pre_handler_kretprobe| that changes the return address when the function
is entered. The return address is set to a trampoline function. On arm, this
function triggers a breakpoiont exception, which eventually calls
\verb|__kretprobe_trampoline_handler|. This function finally calls the user
handler, and sets the exception's return address to the correct one after the
probe.

%\subsection*{Ftrace Tracers}
%TODO
%\subsection*{Fprobe}
%TODO
%\subsection*{Uprobes and Uretprobes}
%TODO
%\subsection*{Eprobes}
%TODO

\section{Perf events}

The perf tracing framework has been described previously. This section will
examine how it is used to register perf events and how they are accessed from
userspace.

Consider the following code of a kernel module registering a perf software
event counter through a PMU:

\begin{Verbatim}[numbers=left, numbersep=5pt, frame=single]
#include <linux/module.h>
#include <linux/kthread.h>
#include <linux/perf_event.h>
#include <linux/uaccess.h>

static struct pmu hello_pmu;
static atomic64_t global_counter = ATOMIC64_INIT(0);

static int hello_pmu_init(struct perf_event *event)
{
  if (event->attr.type != hello_pmu.type)
    return -ENOENT;

  if (is_sampling_event(event))
    return -EINVAL;

  event->hw.state = PERF_HES_STOPPED | PERF_HES_UPTODATE;

  return 0;
}

static void hello_pmu_read(struct perf_event *event)
{
  local64_set(&event->count, atomic64_read(&global_counter));
}

static void hello_pmu_start(struct perf_event *event, int flags)
{
  event->hw.state &= ~(PERF_HES_STOPPED | PERF_HES_UPTODATE);
}

static int hello_pmu_add(struct perf_event *event, int flags)
{
  u64 now = atomic64_read(&global_counter);

  local64_set(&event->count, 0);
  local64_set(&event->hw.prev_count, now);

  if (flags & PERF_EF_START)
    hello_pmu_start(event, flags);

  return 0;
}

static void hello_pmu_stop(struct perf_event *event, int flags)
{
  hello_pmu_read(event);
  event->hw.state |= PERF_HES_STOPPED;
}

static void hello_pmu_del(struct perf_event *event, int flags)
{
  hello_pmu_stop(event, flags);
}

static void hello_pmu_enable(struct pmu *pmu) { }

static void hello_pmu_disable(struct pmu *pmu) { }

static void hello_perf_increment(void)
{
  atomic64_inc(&global_counter);
}

static int __init hello_perf_init(void)
{
  int ret;

  memset(&hello_pmu, 0, sizeof(hello_pmu));

  hello_pmu.task_ctx_nr  = perf_invalid_context;  /* global context */
  hello_pmu.capabilities = PERF_PMU_CAP_NO_INTERRUPT;
  hello_pmu.event_init   = hello_pmu_init;
  hello_pmu.add          = hello_pmu_add;
  hello_pmu.del          = hello_pmu_del;
  hello_pmu.start        = hello_pmu_start;
  hello_pmu.stop         = hello_pmu_stop;
  hello_pmu.read         = hello_pmu_read;
  hello_pmu.pmu_enable   = hello_pmu_enable;
  hello_pmu.pmu_disable  = hello_pmu_disable;

  ret = perf_pmu_register(&hello_pmu, "hello_pmu", -1);
  if (ret)
    return ret;

  pr_info("hello_pmu registered with type %d\n", hello_pmu.type);

  /* Example increment */
  hello_perf_increment();

  return 0;
}

static void __exit hello_perf_cleanup(void)
{
  perf_pmu_unregister(&hello_pmu);
}

module_init(hello_perf_init);
module_exit(hello_perf_cleanup);

MODULE_AUTHOR("Giovanni Santini");
MODULE_DESCRIPTION("perf event example");
MODULE_LICENSE("GPL v2");
\end{Verbatim}

This module defines a counter \verb|global_counter| on line $7$ which is
atomically incremented by the function \verb|hello_perf_increment()|. In the
module initialization function on line $65$, the \verb|hello_pmu| PMU is
initialzied with pointers to functions used to manage it and registered on line
$82$. A notable method is \verb|hello_pmu_read()| on line $22$ which is used to
read the event. This function simply copies the content of the
\verb|global_counter| into the event "count" member. When the module is
unloaded, the PMU is unregistered on line $96$.

An userspace program that reads the event could be implemented with the
following code:

\begin{Verbatim}[numbers=left, numbersep=5pt, frame=single]
#include <linux/perf_event.h>
#include <sys/syscall.h>
#include <unistd.h>
#include <string.h>
#include <stdint.h>
#include <stdio.h>
#include <stdlib.h>

int main(int argc, char **argv)
{
  struct perf_event_attr attr;
  int fd;
  uint64_t count;

  memset(&attr, 0, sizeof(attr));
  attr.type     = 11; /* depends on the specific event */
  attr.config   = 0;
  attr.disabled = 0;
  attr.inherit  = 0;

  /*
  int syscall(SYS_perf_event_open, struct perf_event_attr *attr,
              pid_t pid, int cpu, int group_fd, unsigned long flags);
  */
  fd = syscall(__NR_perf_event_open, &attr, -1, 0, -1, 0);
  if (fd < 0) {
    perror("perf_event_open");
    return EXIT_FAILURE;
  }

  if (read(fd, &count, sizeof(count)) != sizeof(count)) {
    perror("read");
    return 1;
  }

  printf("Current counter value: %lu\n", count);

  close(fd);
  return EXIT_SUCCESS;
}
\end{Verbatim}

This program initializes a \verb|struct perf_event_attr| on line $16$ to $19$
with the information identifying the particular perf event registered in the
module, then calls the \verb|perf_event_open()| systemcall on line $25$ and
reads the counter from the returned file descriptor using the systemcall
\verb|read()| on line $31$.
